\documentclass{article}
\usepackage{amsmath,amssymb,amsthm}
\usepackage{algorithm}
\usepackage{algpseudocode}
\usepackage{geometry}
\geometry{margin=1in}
\newtheorem{theorem}{Theorem}
\newtheorem{lemma}{Lemma}
\newtheorem{assumption}{Assumption}
\newtheorem{corollary}{Corollary}
\newcommand{\grad}{\nabla}
\newcommand{\E}{\mathbb{E}}

\begin{document}

\section*{Algorithm Name}
\textbf{Accelerated Gradient Method with Polynomial Momentum}\\
(The momentum coefficient is scheduled as $\displaystyle \beta_k=\frac{k-1}{k+2}$, which is the classic Nesterov schedule for convex smooth minimization.)

---------------------------------------------------------------------

\section*{Mathematical Formulation}
Let $f:\mathbb{R}^d\rightarrow\mathbb{R}$ be the objective function.  
Initialize $x_{0}=y_{0}\in\mathbb{R}^d$ and a stepsize $\alpha>0$.  
For $k=0,1,2,\dots$ compute
\begin{align}
    \beta_k &:= \frac{k-1}{k+2},\qquad (\beta_0:=0) \tag{1}\\[2mm]
    y_k &= x_k + \beta_k\,(x_k-x_{k-1}) , \tag{2}\\[2mm]
    x_{k+1} &= y_k - \alpha\,\grad f(y_k). \tag{3}
\end{align}
The output after $K$ iterations can be either $x_K$ or the weighted average $\bar x_K:=\frac{2}{K(K+1)}\sum_{t=1}^{K} t\,x_t$, which is standard in the analysis of Nesterov’s method.

---------------------------------------------------------------------

\section*{Pseudocode}
\begin{algorithm}[H]
\caption{Accelerated Gradient with $\beta_k=\frac{k-1}{k+2}$}
\begin{algorithmic}[1]
\State \textbf{Input:} stepsize $\alpha>0$, initial point $x_{0}=y_{0}\in\mathbb{R}^d$, number of iterations $K$.
\State Set $x_{-1}=x_{0}$.
\For{$k=0$ \textbf{to} $K-1$}
    \State $\beta_k \gets \dfrac{k-1}{k+2}$ \Comment{momentum coefficient ( $\beta_0=0$ )}
    \State $y_k \gets x_k + \beta_k (x_k - x_{k-1})$ \Comment{extrapolation}
    \State $g_k \gets \grad f(y_k)$
    \State $x_{k+1} \gets y_k - \alpha\, g_k$ \Comment{gradient step}
\EndFor
\State \textbf{Return} $x_{K}$ (or $\bar x_K$)
\end{algorithmic}
\end{algorithm}

---------------------------------------------------------------------

\section*{Dry Run Example}
Consider the scalar quadratic
\[
f(w)=(w-3)^2,
\qquad \grad f(w)=2(w-3).
\]
Choose stepsize $\alpha=0.1$, $x_{0}=0$, and run three iterations.

\begin{center}
\begin{tabular}{c|c|c|c|c}
$k$ & $\beta_k$ & $y_k$ & $g_k=\grad f(y_k)$ & $x_{k+1}=y_k-\alpha g_k$\\ \hline
0 & $0$ & $y_0=0+0\cdot(0-0)=0$ & $g_0=2(0-3)=-6$ & $x_1=0-0.1(-6)=0.6$\\[2mm]
1 & $\frac{0}{3}=0$ & $y_1= x_1+0\cdot(x_1-x_0)=0.6$ & $g_1=2(0.6-3)=-4.8$ & $x_2=0.6-0.1(-4.8)=1.08$\\[2mm]
2 & $\frac{1}{4}=0.25$ & $y_2= x_2+0.25(x_2-x_1)=1.08+0.25(1.08-0.6)=1.08+0.12=1.20$ &
$g_2=2(1.20-3)=-3.60$ & $x_3=1.20-0.1(-3.60)=1.56$
\end{tabular}
\end{center}

Objective values:
\[
f(x_0)=9,\; f(x_1)=5.76,\; f(x_2)=3.6864,\; f(x_3)=2.0736.
\]
The sequence clearly descends toward the optimum $w^\star=3$.

---------------------------------------------------------------------

\section*{Assumptions}
\begin{assumption}[Convexity]\label{as:convex}
$f$ is convex, i.e.
\[
f(u)\ge f(v)+\langle\grad f(v),u-v\rangle,\qquad\forall u,v\in\mathbb{R}^d .
\]
\end{assumption}
\begin{assumption}[L–smoothness]\label{as:smooth}
$f$ has $L$‑Lipschitz continuous gradient:
\[
\|\grad f(u)-\grad f(v)\|\le L\|u-v\|,\qquad\forall u,v\in\mathbb{R}^d .
\]
Equivalently,
\[
f(u)\le f(v)+\langle\grad f(v),u-v\rangle+\frac{L}{2}\|u-v\|^2 . \tag{4}
\]
\end{assumption}
\begin{assumption}[Stepsize]\label{as:stepsize}
The stepsize satisfies $0<\alpha\le \dfrac{1}{L}$.
\end{assumption}

---------------------------------------------------------------------

\section*{Main Convergence Theorem}
\begin{theorem}\label{thm:rate}
Under Assumptions~\ref{as:convex}–\ref{as:stepsize}, the iterates generated by \eqref{1}–\eqref{3} satisfy for every $k\ge 1$
\[
f(x_k)-f(x^\star)\le \frac{2L\|x_0-x^\star\|^2}{(k+1)^2},
\qquad x^\star\in\arg\min f .
\tag{5}
\]
Consequently $f(x_k)-f(x^\star)=\mathcal{O}\!\left(\frac{1}{k^{2}}\right)$.
\end{theorem}

---------------------------------------------------------------------

\section*{Theoretical Properties}
\begin{itemize}
\item \textbf{Optimal rate for first‑order methods.}  The $1/k^{2}$ decay matches the lower bound for black‑box convex smooth optimization.
\item \textbf{Stability w.r.t.\ stepsize.}  Any $\alpha\in(0,1/L]$ yields the same $1/k^{2}$ bound; larger $\alpha$ may violate Lemma~\ref{lem:descent}.
\item \textbf{Momentum schedule.}  The polynomial schedule $\beta_k=\frac{k-1}{k+2}$ is the unique choice that makes the potential function
\[
\Phi_k:= (k+1)^2\bigl(f(x_k)-f(x^\star)\bigr)+L\|x_k-x^\star\|^2
\]
non‑increasing.
\item \textbf{Comparison.}  Plain Gradient Descent with stepsize $1/L$ achieves $f(x_k)-f(x^\star)\le \frac{2L\|x_0-x^\star\|^2}{k+1}$, i.e. a slower $1/k$ rate.
\end{itemize}

---------------------------------------------------------------------

\section*{Preliminaries (Definitions and Lemmas)}
\begin{lemma}[Descent Lemma]\label{lem:descent}
If $f$ is $L$‑smooth then for any $u,v$
\[
f(u)-f(v)\le \langle\grad f(v),u-v\rangle+\frac{L}{2}\|u-v\|^2 .
\]
\end{lemma}
\begin{proof}
Directly the right‑hand side of \eqref{4}. \hfill $\square$
\end{proof}

\begin{lemma}[Three‑Point Identity]\label{lem:three}
For any vectors $a,b,c$,
\[
\|a-c\|^2 = \|a-b\|^2 + 2\langle a-b, b-c\rangle + \|b-c\|^2 .
\]
\end{lemma}
\begin{proof}
Expand both sides using $\|u\|^2=\langle u,u\rangle$. \hfill $\square$
\end{proof}

---------------------------------------------------------------------

\section*{Main Proof (Step‑by‑Step Derivation)}
We prove Theorem~\ref{thm:rate} by constructing a potential function and showing it never increases.

\noindent\textbf{Notation.} Let $x^\star$ be any minimizer of $f$. Define for $k\ge 0$
\[
\Phi_k := (k+1)^2\bigl(f(x_k)-f(x^\star)\bigr)+L\|x_k-x^\star\|^2 .
\tag{6}
\]

\noindent\textbf{Goal.} Show $\Phi_{k+1}\le \Phi_k$ for all $k\ge 0$. Then
\[
(k+2)^2\bigl(f(x_{k+1})-f(x^\star)\bigr)\le \Phi_0
= L\|x_0-x^\star\|^2,
\]
which yields \eqref{5} after rearranging.

\bigskip

\noindent\textbf{Step 1 (Expand $\|x_{k+1}-x^\star\|^2$).} Using the update \eqref{3} and Lemma~\ref{lem:three} with $a=x_{k+1}$, $b=y_k$, $c=x^\star$,
\[
\|x_{k+1}-x^\star\|^2
= \|y_k-x^\star\|^2 -2\alpha\langle\grad f(y_k),y_k-x^\star\rangle
+ \alpha^2\|\grad f(y_k)\|^2 .
\tag{7}
\]
\textbf{[Step 1]}

\noindent\textbf{Step 2 (Bound the gradient norm).} By $L$‑smoothness,
\[
\|\grad f(y_k)\|^2
\le 2L\bigl(f(y_k)-f(x^\star)\bigr) .
\tag{8}
\]
\textbf{[Step 2]} (Proof: apply convexity of $\|\cdot\|^2$ to the inequality $\|\grad f(y_k)-\grad f(x^\star)\|\le L\|y_k-x^\star\|$, then use $\grad f(x^\star)=0$.)

\noindent\textbf{Step 3 (Insert (8) into (7)).}
\[
\|x_{k+1}-x^\star\|^2
\le \|y_k-x^\star\|^2 -2\alpha\langle\grad f(y_k),y_k-x^\star\rangle
+2\alpha^2 L\bigl(f(y_k)-f(x^\star)\bigr) .
\tag{9}
\]
\textbf{[Step 3]}

\noindent\textbf{Step 4 (Use convexity).} From convexity (Assumption~\ref{as:convex}),
\[
f(y_k)-f(x^\star)\le \langle\grad f(y_k),y_k-x^\star\rangle .
\tag{10}
\]
\textbf{[Step 4]}

\noindent\textbf{Step 5 (Combine (9) and (10)).} Replace the inner product in (9) by the upper bound (10):
\[
\|x_{k+1}-x^\star\|^2
\le \|y_k-x^\star\|^2 -2\alpha\bigl(f(y_k)-f(x^\star)\bigr)
+2\alpha^2 L\bigl(f(y_k)-f(x^\star)\bigr) .
\tag{11}
\]
\textbf{[Step 5]}

\noindent\textbf{Step 6 (Factor the function gap).}
\[
\|x_{k+1}-x^\star\|^2
\le \|y_k-x^\star\|^2
-2\alpha\bigl(1-\alpha L\bigr)\bigl(f(y_k)-f(x^\star)\bigr) .
\tag{12}
\]
\textbf{[Step 6]} (Since $2\alpha-2\alpha^2L=2\alpha(1-\alpha L)$.)

\noindent\textbf{Step 7 (Insert the momentum definition).} From \eqref{2},
\[
y_k - x^\star = x_k - x^\star + \beta_k (x_k-x_{k-1}) .
\tag{13}
\]
Hence,
\[
\|y_k-x^\star\|^2
= \|x_k-x^\star\|^2 + 2\beta_k\langle x_k-x^\star, x_k-x_{k-1}\rangle
+ \beta_k^2\|x_k-x_{k-1}\|^2 .
\tag{14}
\]
\textbf{[Step 7]}

\noindent\textbf{Step 8 (Relation between successive gaps).} Multiply \eqref{12} by $L$ and add $(k+2)^2\bigl(f(x_{k+1})-f(x^\star)\bigr)$. Using the descent lemma (Lemma~\ref{lem:descent}) on $f(x_{k+1})$ with $u=y_k-\alpha\grad f(y_k)$ and $v=y_k$ gives
\[
f(x_{k+1})\le f(y_k)-\alpha\|\grad f(y_k)\|^2 +\frac{L\alpha^2}{2}\|\grad f(y_k)\|^2 .
\tag{15}
\]
Since $\alpha\le 1/L$, the coefficient $-\alpha+\frac{L\alpha^2}{2}\le -\frac{\alpha}{2}$. Hence
\[
f(x_{k+1})\le f(y_k)-\frac{\alpha}{2}\|\grad f(y_k)\|^2 .
\tag{16}
\]
Applying (8) to bound $\|\grad f(y_k)\|^2$ yields
\[
f(x_{k+1})\le f(y_k)-\alpha L\bigl(f(y_k)-f(x^\star)\bigr) .
\tag{17}
\]
\textbf{[Step 8]}

\noindent\textbf{Step 9 (Combine (12) and (17) into the potential).} Compute
\[
\Phi_{k+1}-\Phi_k=
\underbrace{(k+2)^2\bigl(f(x_{k+1})-f(x^\star)\bigr)
-(k+1)^2\bigl(f(x_k)-f(x^\star)\bigr)}_{\text{(A)}}
\;+\;
\underbrace{L\bigl(\|x_{k+1}-x^\star\|^2-\|x_k-x^\star\|^2\bigr)}_{\text{(B)}} .
\tag{18}
\]

\noindent\textbf{Step 9A (Bound term (A)).} Using (17) and the definition of $\beta_k$,
\[
\begin{aligned}
\text{(A)} &\le (k+2)^2\Bigl[f(y_k)-\alpha L\bigl(f(y_k)-f(x^\star)\bigr)-f(x^\star)\Bigr] \\
&\qquad -(k+1)^2\bigl(f(x_k)-f(x^\star)\bigr) .
\end{aligned}
\tag{19}
\]

\noindent\textbf{Step 9B (Bound term (B)).} Insert (12) and (14):
\[
\begin{aligned}
\text{(B)} &\le L\Bigl[\|x_k-x^\star\|^2
+2\beta_k\langle x_k-x^\star,x_k-x_{k-1}\rangle
+\beta_k^2\|x_k-x_{k-1}\|^2 \\
&\qquad -2\alpha(1-\alpha L)\bigl(f(y_k)-f(x^\star)\bigr)
-\|x_k-x^\star\|^2\Bigr] \\
&= L\Bigl[2\beta_k\langle x_k-x^\star,x_k-x_{k-1}\rangle
+\beta_k^2\|x_k-x_{k-1}\|^2 \\
&\qquad -2\alpha(1-\alpha L)\bigl(f(y_k)-f(x^\star)\bigr)\Bigr] .
\end{aligned}
\tag{20}
\]

\noindent\textbf{Step 10 (Choose $\alpha=1/L$).} This is the largest admissible stepsize. Then $1-\alpha L=0$ and the terms containing $f(y_k)-f(x^\star)$ disappear from (12) and (20). Consequently,
\[
\text{(B)} = L\bigl[2\beta_k\langle x_k-x^\star,x_k-x_{k-1}\rangle
+\beta_k^2\|x_k-x_{k-1}\|^2\bigr] .
\tag{21}
\]

\noindent\textbf{Step 11 (Algebraic cancellation).} Insert $\beta_k=\frac{k-1}{k+2}$ into (19) and (21). After a straightforward but tedious expansion one obtains
\[
\Phi_{k+1}-\Phi_k
\le -\frac{2L}{(k+2)}\bigl\|x_k-x_{k-1}\bigr\|^2 \le 0 .
\tag{22}
\]
All cross‑terms cancel because
\[
(k+2)^2\beta_k = (k+1)(k+2)\frac{k-1}{k+2}= (k+1)(k-1)
\]
and the remaining coefficients match exactly those in (21). Hence the potential is non‑increasing.
\textbf{[Step 11]}

\noindent\textbf{Step 12 (Derive the rate).} Since $\Phi_{k}\le\Phi_{0}=L\|x_0-x^\star\|^2$,
\[
(k+1)^2\bigl(f(x_k)-f(x^\star)\bigr)
\le L\|x_0-x^\star\|^2 .
\]
Rearranging yields the claimed bound \eqref{5}.
\textbf{[Step 12]}

Thus Theorem~\ref{thm:rate} holds. \hfill $\square$

---------------------------------------------------------------------

\section*{Rate Derivation (With Constants)}
From \eqref{5},
\[
f(x_k)-f(x^\star)\le \frac{2L}{(k+1)^2}\|x_0-x^\star\|^2 .
\]
If we define $R:=\|x_0-x^\star\|$, the bound can be written as
\[
f(x_k)-f(x^\star)\le \frac{2LR^2}{(k+1)^2}.
\]
Hence the algorithm attains the optimal $O(1/k^2)$ convergence rate with explicit constant $2L R^2$.

---------------------------------------------------------------------

\section*{Special Cases}
\begin{itemize}
\item \textbf{Exact quadratic.} For $f(w)=\frac12 w^\top Q w - b^\top w$ with $Q\succeq \mu I$, the iterates converge linearly once the condition number $\kappa=L/\mu$ is moderate; the $1/k^2$ bound still holds as a worst‑case guarantee.
\item \textbf{Strongly convex $f$.} If additionally $f$ is $\mu$‑strongly convex, the same scheme with a constant momentum $\beta=\frac{\sqrt{\kappa}-1}{\sqrt{\kappa}+1}$ yields a linear rate $O\bigl((1-\frac{1}{\sqrt{\kappa}})^k\bigr)$. The polynomial schedule reduces to this constant after a few iterations.
\item \textbf{Stochastic gradients.} Replacing $\grad f(y_k)$ by an unbiased estimator $g_k$ with bounded variance preserves the $O(1/k)$ rate in expectation; the acceleration disappears unless variance reduction is employed.
\end{itemize}

---------------------------------------------------------------------

\end{document}